\begin{abstract}
Feed-forward dynamic 3D Gaussian Splatting (3DGS) methods have recently
made it possible to reconstruct dynamic scenes from uncalibrated monocular
video in near real time. However, because they operate in a normalized
canonical space with an orthographic camera, the resulting
representations are \emph{scale-ambiguous}: a reconstructed scene cannot be
used for metric reasoning in robotics, autonomous driving, or augmented
reality without a separate, error-prone scale-recovery step. We present
\StereoSplat, a fully feed-forward framework that reconstructs
\emph{metric-scale} dynamic 3D Gaussian scenes from a calibrated stereo
video stream. The key insight is that the stereo baseline and focal length
provide a closed-form physical anchor for absolute scale:
$Z = b f / d$. Building on the \StreamSplat architecture, we replace the
monocular depth encoder with a stereo-matching front end that produces
metric depth, swap the orthographic rasterizer for a perspective one that
respects the true intrinsics and per-frame poses, and introduce a
\emph{metric back-projection head} that lifts per-pixel Gaussian
attributes into the metric world frame. To address the well-documented
display artifacts of dynamic GS --- temporal flickering, popping, and
high-frequency motion blur --- we further integrate a streaming display
module with a motion-adaptive temporal smoother and an opacity-stability
filter. Experiments on TartanAir, DSEC, and KITTI Stereo show that
\StereoSplat recovers metric scale to within \TODO{X}\,m L1 error while
matching or surpassing the photometric quality of monocular baselines, and
that each component of the proposed pipeline contributes measurably to the
final result.
\end{abstract}
